\documentclass[10pt,a4paper]{amsart}
\usepackage[margin=19mm]{geometry}
\usepackage{amsmath,amssymb,mathtools}
\usepackage[scr=rsfs,cal=euler]{mathalfa}
\usepackage{xfrac}
\usepackage[unicode,hidelinks,bookmarks=false]{hyperref}
\newcommand{\ie}{i.e.\ }
\newcommand{\cf}{cf.\ }
\newcommand{\resp}{respectively\ }
\newcommand{\Subsection}{\subsection}
\providecommand{\nobreakdash}{\nobreak\hbox{-}}
\setlength{\emergencystretch}{2em}
\multlinegap0pt
\usepackage{amssymb}
\usepackage{mathtools}
\usepackage[shortlabels]{enumitem}
\usepackage{xcolor}
\usepackage[normalem]{ulem}
\usepackage{float}
\newtheorem{proposition}{Proposition}
\DeclareMathOperator{\CR}{CR}
\DeclareMathOperator{\inrad}{inrad}
\newcommand{\norm}[1]{\lVert #1\rVert}
\title{\textbf{Bellman Search in Arbitrary Finite Dimension:}\\ A Self-Similar Cell Theorem and Effective Computability of Planar Shoreline Search}
\author{Florentin Koch}
\date{Source: arXiv:2608.29032v1 (CC BY 4.0)}
\begin{document}
\maketitle

\noindent\emph{Source and licence.} \textbf{Bellman Search in Arbitrary Finite Dimension:}\\ A Self-Similar Cell Theorem and Effective Computability of Planar Shoreline Search. Version arXiv:2608.29032v1 (CC BY 4.0), distributed by arXiv under the Creative Commons Attribution 4.0 International licence, http://creativecommons.org/licenses/by/4.0/, as recorded in the arXiv licence metadata for this exact version and shown on the abstract page https://arxiv.org/abs/2608.29032v1. Full licence notice: \texttt{licenses/arxiv-2608.29032-CC-BY-4.0.txt}.\par\smallskip

\noindent\textit{Editorial scope.} Counted unit: the article's proposition prop:compact-bounds (source0.tex lines 853--865). The article's complete original proof of it is delivered with it (source0.tex lines 867--895, one \texttt{proof} environment of 29 lines). No mathematics is written, repaired or re-derived: the statement and the proof are the article's own lines, in the article's own notation, and no retained line is substituted or re-flowed.  Retained uncounted as the source-local context the proof needs: lines 822--825.  Nothing was omitted from any retained span, so the package carries no omission record.  No retained line contains a citation, so the wrapper carries no bibliography and no bibliography entry.  No retained line contains a figure environment, an image-inclusion command or drawing code, so no image is delivered and the package declares no asset dependency.  Editorial formatting only, in addition: an AMS \texttt{amsart} preamble that restates the article's own declarations and macros used by the retained lines, namely the environments \texttt{proposition} on separate counters, together with the standard packages those lines need.  The retained environments print in this document as Proposition 1 in order of appearance, whereas the article prints the same items with the numbering of its own full document.  The complete native source of the article as distributed in the arXiv e-print for this exact version is bundled unchanged as \texttt{sources/208823/source0.tex}.
\begin{equation}
\inrad_0(P)=1,
\label{eq:inradius-normalization}
\end{equation}

\begin{proposition}[Compact bounds]
\label{prop:compact-bounds}
Fix $N$ and $U>2$. Every polygonal cell with $N$ nonzero edges satisfying \eqref{eq:inradius-normalization} and $\CR(p,Q)\le U$ obeys
\begin{equation}
2\le\ell\le U,
\qquad
\max_i\norm{p_i}\le U-1,
\qquad
\frac{1}{U^N-1}\le\lambda\le1-\frac{2}{U}.
\label{eq:compact-bounds}
\end{equation}
The corresponding sublevel is therefore compact after zero-length subdivisions are allowed, and its infimum is attained.
\end{proposition}

\begin{proof}
At the beginning of the cell, the past has inradius $\lambda$ and cost $\lambda\ell/(1-\lambda)$, hence $\ell/(1-\lambda)\le U$. Since $B_2\subset P$, the cell joins two points separated by at least $2$, so $\ell\ge2$ and $\lambda\le1-2/U$. If $R:=\max_i\norm{p_i}$, a direction opposite a farthest vertex provides a point of the hull with support at most $-1$; thus $\ell\ge R+1$ and $R\le U-1$.

For the lower bound on $\lambda$, write $S_e:=\ell_1+\cdots+\ell_e$ and choose, at the end of edge $e$, a unit direction $u$ opposite to $p_e-p_{e-1}$. Since $p_N=Qp_0$ and $\lambda=Q^{-1}$, one has $p_0=\lambda p_N\in\lambda P$. For every previously visited vertex $p_j$, $j<e$,
\[
u\cdot p_j\le u\cdot p_0+\norm{p_j-p_0}
\le h_{\lambda P}(u)+S_{e-1}
\le\lambda R+S_{e-1}.
\]
The new point does not increase support in direction $u$. Hence the available support at the end of edge $e$ is at most $\lambda R+S_{e-1}$, and the ratio constraint gives
\[
S_e\le U(\lambda R+S_{e-1}).
\]
Iteration yields
\[
S_e\le\lambda R U\frac{U^e-1}{U-1}.
\]
At $e=N$, with $S_N=\ell$, $R\le\ell-1$, and $\ell\le U$,
\[
\lambda\ge
\frac{\ell(U-1)}{RU(U^N-1)}
\ge
\frac{\ell(U-1)}{(\ell-1)U(U^N-1)}
\ge
\frac{1}{U^N-1},
\]
where the last inequality uses $\ell/(\ell-1)\ge U/(U-1)$ for $2\le\ell\le U$.
This proves \eqref{eq:compact-bounds}.
\end{proof}

\end{document}
